\chapter{Parallel RLC Circuits}
\label{ch:rlcpar}
The parallel RLC circuit is the \emph{dual} of the series circuit of
\cref{ch:rlc}: same two energy stores, same second-order dynamics, same
\(\omega_0\)---but the roles of voltage and current, impedance and admittance, are
exchanged. That duality does most of the work, so this chapter is short. It states
the model, then concentrates on the two things that genuinely differ: the circuit
looks like a high impedance at resonance instead of a low one, and its quality
factor is built the other way up.

The parallel form matters in radio because it is the \textbf{tank circuit}---the
resonator in an oscillator, the tuned load of a Class C amplifier, the trap that
blocks one frequency. Wherever a circuit needs to present a high impedance at one
frequency and a low impedance elsewhere, it is a parallel resonator.

\section{The Circuit}
Drive \(R\), \(L\), and \(C\) in parallel from a current source
(\cref{fig:rlcpar-schematic}). The dual of ``one loop, one current'' is ``one node,
one voltage'': all three elements share the node voltage \(v\), which we take as the
output, \(v_o=v\). The two states are still one per store---the capacitor voltage
(here the same as \(v\)) and the inductor current \(i_L\).

\begin{figure}[htbp]
\centering
\begin{circuitikz}[scale=1]
\draw (0,0) to[I,l=\(u(t)\)] (0,2.4);
\draw (0,2.4) -- (5.4,2.4);
\draw (0,0) -- (5.4,0);
\draw (1.6,2.4) to[R,l=\(R\)] (1.6,0);
\draw (3.2,2.4) to[L,l=\(L\),i>=\(i_L\)] (3.2,0);
\draw (4.8,2.4) to[C,l=\(C\)] (4.8,0);
\draw (5.4,2.4) to[short,-o] (6.3,2.4) node[right]{\(+\)};
\draw (5.4,0) to[short,-o] (6.3,0) node[right]{\(-\)};
\node at (6.7,1.2) {\(v_o\)};
\node[ground] at (2.4,0) {};
\end{circuitikz}
\caption{Parallel RLC circuit driven by a current source. Every element shares the
node voltage, which is the output \(v_o=v\). The states are the capacitor voltage
\(v\) and the inductor current \(i_L\).}
\label{fig:rlcpar-schematic}
\end{figure}

\section{From Schematic to Differential Equation}
Where the series circuit used KVL around a loop, the parallel circuit uses KCL at
the node: the source current splits among the three branches,
\[
u(t)=\underbrace{\frac{v}{R}}_{\text{resistor}}
+\underbrace{i_L}_{\text{inductor}}
+\underbrace{C\frac{\dd v}{\dd t}}_{\text{capacitor}},
\qquad
v=L\frac{\dd i_L}{\dd t}.
\]
As before this is a coupled pair of first-order equations, one per store:
\[
\frac{\dd v}{\dd t}=\frac{1}{C}\Bigl(u-\frac{v}{R}-i_L\Bigr),
\qquad
\frac{\dd i_L}{\dd t}=\frac{1}{L}\,v .
\]
Eliminating \(i_L\) (differentiate the first, substitute the second) gives the
second-order equation in the node voltage:
\[
\boxed{\,LC\frac{\dd^2 v}{\dd t^2}+\frac{L}{R}\frac{\dd v}{\dd t}+v
=L\frac{\dd u}{\dd t}.\,}
\]
Compare this with the series result \(LC\ddot v_C+RC\dot v_C+v_C=u\). The
second-derivative term is identical, so \(\omega_0\) will be identical; only the
damping term differs, and it differs by having \(R\) in the \emph{denominator}.
That single swap is the whole story of this chapter.

\section{State-Space Form and the Generic Match}
With \(x_1=v\) and \(x_2=i_L\),
\[
\dot{\state}=
\mat{-1/(RC) & -1/C\\ 1/L & 0}\state
+\mat{1/C\\ 0}u,
\qquad
y=v_o=\mat{1 & 0}\state .
\]
The characteristic polynomial is
\[
\det(sI-A)=s\Bigl(s+\frac{1}{RC}\Bigr)+\frac{1}{LC}
=s^2+\frac{1}{RC}s+\frac{1}{LC},
\]
so matching the standard form \(s^2+2\zeta\omega_0 s+\omega_0^2\) gives
\[
\omega_0=\frac{1}{\sqrt{LC}},
\qquad
\alpha=\frac{1}{2RC},
\qquad
\zeta=\frac{\alpha}{\omega_0}=\frac{1}{2R}\sqrt{\frac{L}{C}},
\qquad
Q_p=\frac{1}{2\zeta}=R\sqrt{\frac{C}{L}}=\frac{R}{\omega_0 L}=\omega_0 RC .
\]
The natural frequency is the same \(1/\sqrt{LC}\); the damping is inverted.

\section{Parallel Resonance: a High-Impedance Peak}
Work in admittances, which is what parallel networks add (\cref{ch:ac}):
\[
Y(\jj\omega)=\frac{1}{R}+\jj\omega C+\frac{1}{\jj\omega L}
=\frac{1}{R}+\jj\left(\omega C-\frac{1}{\omega L}\right).
\]
At \(\omega=\omega_0=1/\sqrt{LC}\) the two susceptances cancel exactly---the
capacitor's \(+\omega C\) against the inductor's \(-1/(\omega L)\)---leaving
\(Y=1/R\), purely real and \emph{minimum}. Since \(Z=1/Y\), the impedance is
therefore purely resistive and \emph{maximum}:
\[
Z(\jj\omega_0)=R .
\]
In the lossless ideal (\(R\to\infty\)) the input impedance becomes infinite: the
tank is an open circuit at its resonance. Real loss keeps it finite.

This is the exact mirror of the series case, and the contrast is worth stating
plainly:
\begin{center}
\begin{tabularx}{\textwidth}{@{}L{0.26\textwidth}Y Y@{}}
\toprule
At resonance & Series RLC & Parallel RLC \\
\midrule
Reactances/susceptances & \(X_L\) cancels \(X_C\) & \(B_C\) cancels \(B_L\) \\
Impedance \(Z(\jj\omega_0)\) & \(R\) --- a \emph{minimum} & \(R\) --- a \emph{maximum} \\
Current from the source & maximum & minimum \\
Behaves like & a short at \(\omega_0\) & an open at \(\omega_0\) \\
Internal magnification & \(v_L,v_C\) exceed source by \(Q\) & circulating tank current exceeds source by \(Q\) \\
Quality factor & \(Q_s=\omega_0 L/R\) & \(Q_p=R/(\omega_0 L)\) \\
\bottomrule
\end{tabularx}
\end{center}

\begin{physicalbox}
\textbf{Why the circulating current is large.} At resonance the source only has to
replace what \(R\) dissipates, so the current it supplies is small. But \(L\) and
\(C\) are meanwhile trading a large current back and forth between themselves---
\(Q\) times the source current---because each one's reactive current is cancelled by
the other's as far as the source can see. The series circuit does the dual thing
with voltage: a modest source voltage while \(v_L\) and \(v_C\) individually swing
\(Q\) times larger, equal and opposite.
\end{physicalbox}

\section{Why \texorpdfstring{\(Q\)}{Q} Inverts}
The two quality factors look contradictory until you ask where the loss sits:
\[
Q_s=\frac{\omega_0 L}{R}
\qquad\text{versus}\qquad
Q_p=\frac{R}{\omega_0 L} .
\]
In the series circuit the loss resistance is \emph{in the path} of the resonating
current, so a \emph{larger} \(R\) dissipates more per cycle and lowers \(Q\). In
the parallel circuit the loss resistance is a \emph{shunt across} the resonator, so
a larger \(R\) bleeds off \emph{less} current and \emph{raises} \(Q\). A perfect
series resonator wants \(R\to0\); a perfect parallel resonator wants
\(R\to\infty\). Both statements say the same physical thing---minimize the energy
lost per cycle---which is the definition \(Q\) actually encodes:
\[
Q=2\pi\,\frac{\text{energy stored}}{\text{energy lost per cycle}} .
\]
Read that way there is only one \(Q\), and the two formulas are what it becomes in
two topologies.

\begin{mistakebox}
Do not carry a \(Q\) formula across topologies. Ask first whether the loss is in
series with the resonating current or in parallel with the resonator, then use the
matching form. (The loaded/unloaded/component distinction of \cref{ch:rlc} applies
here unchanged.)
\end{mistakebox}

\begin{workedbox}[: the circulating current is real too]
Take the dual of \cref{ch:rlc}'s magnification box: \(L=\SI{20}{\micro\henry}\),
\(C=\SI{100}{\pico\farad}\), so \(f_0\approx\SI{3.56}{\mega\hertz}\) and
\(\omega_0L\approx\SI{447}{\ohm}\) as before---but now the loss is a
\(R=\SI{100}{\kilo\ohm}\) shunt, giving
\[
Q_p=\frac{R}{\omega_0L}=\frac{100000}{447}\approx224 .
\]
Drive it with \SI{10}{\milli\ampere} RMS from a current source. At resonance
\(Z=R=\SI{100}{\kilo\ohm}\), so the node develops
\(V=(0.010)(100000)=\SI{1}{\kilo\volt}\), and the current circulating in the
inductor is
\[
I_L=\frac{V}{\omega_0L}=\frac{1000}{447}\approx\SI{2.2}{\ampere},
\]
which is \(Q\) times the \SI{10}{\milli\ampere} the source supplies. Compare the
series box: there a modest voltage produced \(Q\)-fold \emph{voltages} across the
elements; here a modest current produces a \(Q\)-fold \emph{current} around the
loop. Same \(Q\), same \(f_0\), dual variables---and the same practical warning,
since that circulating current is what heats a tank coil.
\end{workedbox}

The bandwidth statement, however, is topology independent:
\[
BW\approx\frac{f_0}{Q},
\]
because it is a statement about pole geometry---the distance from the
\(\jj\omega\) axis---not about wiring (\cref{ch:splane}).

\section{Root Locus and Energy Balance}
Both of these carry over from \cref{ch:rlc} by duality, so they need only a sentence
each.

The \emph{root locus} is identical in shape: the poles are
\(s=-\alpha\pm\jj\omega_d\) with the same \(\omega_0\), so sweeping the loss
resistance still walks them along the circle of radius \(\omega_0\) from the
\(\jj\omega\) axis to the real axis (\cref{fig:rlc-rootlocus}). Only the sweep
direction reverses: for the series circuit the poles move away from the axis as
\(R\) \emph{increases}, while here they move away as \(R\) \emph{decreases}. For the
parallel circuit the more informative plot is usually the \emph{impedance} locus,
which is what \cref{fig:rlcpar-worked} shows.

The \emph{energy balance} is the dual statement of the series one. With
\(\mathcal{H}=\tfrac12 Cv^2+\tfrac12 Li_L^2\),
\[
\frac{\dd\mathcal{H}}{\dd t}=u(t)\,v(t)-\frac{v^2(t)}{R},
\]
where the loss term is now \(v^2/R\) instead of \(Ri^2\)---the same physics through
the dual variable.

\begin{energybox}
High \(Q_p\) again means little energy lost per cycle, so the \(L\)--\(C\) exchange
rings for many cycles. Note that the \emph{energy} definition of \(Q\) is identical
in both topologies even though the two formulas look reciprocal: it is always
stored energy over energy lost per cycle, and only the location of the loss changes.
\end{energybox}

\section{Worked Example}
Take the same \(L=\SI{10}{\micro\henry}\) and \(C=\SI{220}{\pico\farad}\) as the
series example of \cref{ch:rlc}, but now in parallel with a loss resistance
\(R=\SI{18}{\kilo\ohm}\) (a realistic unloaded tank). Then
\[
f_0=\frac{1}{2\pi\sqrt{LC}}\approx\SI{3.39}{\mega\hertz}
\qquad\text{(unchanged---same \(L\) and \(C\))},
\]
\[
\omega_0 L=(2.13\times10^{7})(10\times10^{-6})\approx\SI{213}{\ohm},
\qquad
Q_p=\frac{R}{\omega_0 L}=\frac{18000}{213}\approx 84 .
\]
The half-power bandwidth is \(BW\approx f_0/Q_p\approx\SI{40}{\kilo\hertz}\), and at
resonance the tank presents \(Z=\SI{18}{\kilo\ohm}\), purely resistive---a high
impedance an oscillator or Class C stage can develop signal voltage across. Note
the contrast with \cref{ch:rlc}: there, \(R=\SI{5}{\ohm}\) in series gave
\(Q_s\approx43\); here \(R=\SI{18}{\kilo\ohm}\) in parallel gives \(Q_p\approx84\).
Small series \(R\) and large parallel \(R\) both mean a good resonator.

\Cref{fig:rlcpar-worked} plots it. Set this against the series example's
\cref{fig:rlc-phasor} and the duality becomes visual rather than verbal: the series
circuit's \(|Z|\) \emph{dips} to \(R\) at \(f_0\) while this one \emph{peaks} at
\(R\); the series impedance locus is a straight vertical line while this one is a
circle through the origin. Both are the same second-order pole pair, read through
reciprocal coordinates.

\begin{figure}[htbp]
\centering
\includegraphics[width=\textwidth]{rlc_parallel_worked}
\caption{The parallel worked example, \(L=\SI{10}{\micro\henry}\),
\(C=\SI{220}{\pico\farad}\), \(R=\SI{18}{\kilo\ohm}\). \emph{Left:} the tank
impedance magnitude peaks at \(|Z|=R=\SI{18}{\kilo\ohm}\) at
\(f_0=\SI{3.39}{\mega\hertz}\), falling away on both sides with a half-power
bandwidth of about \SI{40}{\kilo\hertz}---the opposite of the series circuit's dip.
\emph{Right:} the impedance locus is a circle through the origin, inductive below
\(f_0\) and capacitive above (again the reverse of the series case), crossing the
real axis at \(Z=R\).}
\label{fig:rlcpar-worked}
\end{figure}

\begin{exambox}
Exam questions on resonance almost always hinge on one decision: series or
parallel? Series resonance means minimum impedance and maximum current; parallel
resonance means maximum impedance and minimum \emph{source} current with maximum
\emph{circulating} current. Both share \(f_0=1/(2\pi\sqrt{LC})\), and in both
\(R\) sets \(Q\) and bandwidth but never \(f_0\).
\end{exambox}
